\emph{Carried over verbatim from the full report (draft v2.4, 23 September 2026), which is released with this report. Section and appendix numbers in this appendix refer to that report; Appendix~E maps them to this one.}

\emph{Its paragraph on the use of AI assistants is superseded by the AI usage statement of this report, which records a second exception (submission 1016734).}

\subsection*{Full report, Section 12: Limitations}

Ground truth under the earlier labeling rule exists for two evaluated models: 164 pairs for DeepSeek, 26 of whose labels we derived (2.2), and 10 for gpt-oss-120b at low reasoning effort, five of them with a declared label. Every probe and every concordance in Sections 7 and 8 is measured on DeepSeek and assumed, not shown, for the other four. The main-phase train split (2.1) labels all five evaluated models under the new rule, but with only 3 to 37 labeled rows per model, and apart from the evaluation of our entry and the five-row comparison in 6.4 the analyses in this draft do not use it. Sections 8.1 and 8.2 use 40 and 174 problems and one model. The fourteen-state structure is established; its cause is inferred, not proven: a choice of kernels or algorithms made when the process starts, which either disabling compilation of the static-cache forward or forcing deterministic algorithms removed in our sequential runs. The remedies were tested only in sequential runs, each on one card. The simulated 0.94 in 7.3 treats observed frequencies as true probabilities, and on our own labels the same procedure overstates the measured reliability by about 0.4. The regeneration in 7.2 used gpt-oss-120b prompted with the public examples and gated by itself; the public records attribute the organizers' own variants of the same families to language models as well, so the negative concordance is a fact about our generation, gating, answer key, sampling counts, decoding and prompting, not about the families.

Several analyses were first run from scripts that were not saved in the project. On 23 September three were re-created and saved: the subsampling curve of 7.2 and the second-level length control of 8.3, both extended from one draw or split to twenty (the second gives a gain of about two points, against the 0.8 of the single split first reported), and the simulation of 7.3, re-created from its dated description. The same day, scripts were written and saved, with their outputs, for the cross-run analysis of E14 (8.1), the AUROCs, rates and accuracies of E16 (2.3b, 8.2, 8.4), the controls of the public data (Section 2, 6.3), the artifact controls and the nested null of 8.3, the recomputations of the warmup hypotheses (Section 4), the permutation p-values of E11 and the other controls the second and third verifications required (Section 13). A few figures still rest only on a dated record: the ICC values of H29, written to a Modal volume that was not downloaded; facts read from container logs at the time and recorded in dated files (the item counts per model of 3.2 and 6.1, the 119 s per item of 907677, V8's per-item predictions on the current tasks, v6's runtimes and gpt-oss thresholds, and per-model score ranges); Codabench scores and leaderboard facts recorded in dated files and absent from our saved snapshots (the warmup scores of 907310, 907677, 909843, 909853, 909994, 910002 and 910013, V8's scores of 13 and 14 September on the original tasks, and the leaderboard figures of 29 August and of 13 and 14 September); and the compute totals of Section 13. Nothing here estimates performance on the final held-out problems; Appendix C states what we predicted before knowing.

\subsection*{Full report, Section 13: Reproducibility and disclosure}

\paragraph{Pre-registrations.}
PREREG\_E1 to E4, E6 to E10, E12 to E15 and E13b, and V4 to V8b, each dated before execution. The hypothesis tests among them carry kill criteria; E13 (a reliability measurement read through a pre-registered grid), E14 and V7 (measurements with written predictions) carry none, and V4, V8 and V8b state predictions and decision rules. Declared deviations (E6, E12, E13, E14, E15, and V4, whose entry was submitted once as a measurement although its pre-registration said a refuted entry would not be sent, Appendix B) are preserved in place, and dated correction notes were appended on 21 and 23 September. E5 is documented only in its outcome file; E11 ran from a written plan (PIANO\_E11\_definitivo) without the frozen pre-registration it called for; E16 (8.2, 8.4) was not pre-registered, although our reply of 7 September to the organizers said every experiment on their VM would be. Deviations declared here for the first time: E4 (H7 compared with the global prior of the training models instead of the pre-registered per-model majority, against which it scores $-$1.2 points); E6 (H10 run as a logistic regression on 13 of the 14 features in place of the pre-registered regressor and threshold); E11 (three planned primary tests not run; a Holm correction planned over five tests at alpha 0.010, announced over two and not applied; the confirmation set not held out; AUROCs computed on the raw scalars with the orientation chosen on the data; the text baseline reduced to character count; the stratified and bookkeeping reference AUROCs not reported; only the primary depth extracted); E13 (gpt-oss-120b not labeled; DeepSeek's A$'$ condition labeled with a different prompt; the per-family subsampling curve replaced by a per-problem grid, 7.2; the pre-registered reading grid not followed, 7.3); E14 (measure m4 not computed); and E15 (in the nested null, 60 permutations instead of 500 and two of the four regularization values, the latter re-run with the full grid on 23 September; Holm over the two primaries not computed).

\paragraph{What we logged, and what we know about the evaluation sets.}
Container measurements: submissions 907310, 907677, 907967, 909567, 909843, 909853, 909976, 909994, 910002, 910013, 912096, 912097, 924968, 924972, 925822, 936802, 936806. They logged aggregate counts, the number of items per model, per model batch the minimum and maximum regressor score and the threshold applied (v5 to v7), per-model length summaries (V8: minimum, median and maximum character count, which on batches of three items or fewer are every item's length, unattributed), timings, model identifiers and, in 907967, the GPU, memory and library versions and the prompt and generated token counts of one DeepSeek warmup item; we read those logs. No problem text and no per-item label from the evaluation sets was logged, stored or recovered. Aggregate facts about those sets were observed or inferred and are reported where used: item counts per model (3.2, 6.1), the class balance implied by the pre-registered null probe and by constant entries (3.1, 6.3) and per-model score ranges (8.4). One warmup submission was built on aggregate facts about the evaluation set and is disclosed in Appendix B. Logged aggregates also shaped two other design choices, neither of which used evaluation labels: H17 (v6) calibrated a separate threshold for gpt-oss-120b because its logged scores on the warmup set (9 of 17 items, 0.340 to 0.870) straddled the inherited threshold, while the threshold itself was fitted on the ten gpt-oss rows of \texttt{val-sample}, one of them robust (6.2); and the budget guard's per-item cost estimate was set from container timings (130 s after the measured 119 s, then 90 s after the 21 to 59 s logged by v5).

\paragraph{Compute.}
GPU time on Modal: the dated files document about 12 to 13 GPU-hours and about 34 dollars, of which 9.2 A100-hours (about 19 dollars, funded personally) went to the extraction of Section 5 and the rest to its pilot and to E9 to E12; the billed total, including start-up and loading time and how much of it platform credit covered, will be itemized from the billing records before release. On the organizers' VM, E13 took about 100 card-hours by the estimate in its outcome file, and E14 and E16 several more; a total will be itemized from the run logs before release.

\paragraph{Release.}
Code, extraction outputs (hashes, scores, per-variant answers), the analysis scripts and every dated file will be released with the final version, including drafts v1.0 and v2.1 to v2.4 (draft v2.0, of 14 September, was edited in place into v2.1 and was not kept as a separate file; its four corrections are items 1 to 4 of Appendix D). In the dated files, the names, handles and personal details of other participants, and every passage about their scores, methods or conduct, will be redacted, with each redaction marked, as in Appendix C. The invalidated Qwen labeling runs are kept with a file stating why. The analysis scripts saved on 21 and 23 September are \texttt{controllo\_h2\_alias\_2026-09-21.py}, in the project root (H2 with the duplicated rows counted once, Section 4), and, in \texttt{vm/}: \texttt{e13b\_regola\_nuova.py}, \texttt{e13b\_controllo\_modello.py}, \texttt{e15\_percorso\_ufficiale.py}, \texttt{simulazione\_affidabilita\_organizzatori.py}, \texttt{e13\_sottocampionamento.py}, \texttt{e15\_controllo\_lunghezza.py}, \texttt{e15\_null\_nidificato.py}, \texttt{e14\_tabella.py}, \texttt{e16\_analisi.py}, \texttt{e15\_artefatto\_controlli.py}, \texttt{e11\_holm.py}, \texttt{warmup\_ricalcoli.py}, \texttt{dati\_pubblici\_controlli.py}, \texttt{controlli\_verifica2.py}, \texttt{controlli\_verifica3.py}, \texttt{controlli\_verifica4.py} and \texttt{v8\_train\_main\_v2.py}, with their outputs in \texttt{vm\_risultati/}; the Codabench snapshots of 7, 21 and 23 September and the three files downloaded from the organizers' baselines repository (\texttt{prompts/solve.txt}, \texttt{data/math-robust-agg.csv} and \texttt{scripts/probe.py}) are in \texttt{dati/}.

\paragraph{Reproducing our score.}
Our entry on the current tasks is V8. The Main bundle \texttt{solution-v8-main.zip} (SHA-256 \texttt{7b70f15c51752dcf4285e477d58c541d0ecf052d833c325c0702dbbd054ffb5f}) contains a single \texttt{solution.py} that defines \texttt{are\_robust(model\_id, problems)}, returns \texttt{[len(p) >= 98 for p in problems]} and logs only per-model counts and length summaries; the Small bundle \texttt{solution-v8-small.zip} (SHA-256 \texttt{0b83b55f94dc16b9852a7c2b35d71ac5787e43d81afd6396dee00fb5e011e500}) adds the empty marker file \texttt{small.txt}. Upload each bundle with only its own task selected: Main on task 36237 (expected 10/21, submission 936802), Small on task 36241 (expected 7/14, submission 936806). The rule loads no model and reads no data at run time. Its two parameters were fixed on the six public \texttt{aimo-interp} datasets, which are its training data: the threshold, 98, is the lower median of the character counts of their 174 distinct problems (the median is 98.5), and the direction, longer problems robust, is the one with the higher balanced accuracy on their labeled pairs (0.652 over the 335 published pairs, against 0.348 for the reverse). \texttt{vm/controlli\_verifica3.py} recomputes the threshold and \texttt{vm/dati\_pubblici\_controlli.py} the balanced accuracies.

\paragraph{Use of AI assistants.}
This work was carried out with an AI assistant (Claude, by Anthropic, used through Claude Code). It drafted the pre-registrations, including their hypotheses, predictions and kill criteria; wrote the solution code of our submission bundles and the analysis and extraction scripts, and ran them; operated the organizers' VM over SSH for E13 to E16; made most of our Codabench uploads through the author's browser, at the author's instruction (recorded in PREREG\_V4 for 909843 and in PREREG\_V8b for 936802 and 936806), except the resubmission 925822 of 14 September, which it made without a specific instruction, to restore what it had wrongly taken to be a lost Main row (6.3); read the public Codabench, GitHub, Hugging Face and competition-website records, the organizers' Discord channels, their proposal on arXiv, the works cited in Section 1 (on arXiv, the NeurIPS website and the Thinking Machines Lab website) and the author's email correspondence with them, behind the dates, quotations and leaderboard figures of the header, of Sections 1, 2, 3, 6, 7, 8, 9, 10 and 13 and of Appendices B to D; drafted the correspondence and this report; and ran automated multi-agent adversarial reviews of hypotheses, drafts and code, including the verifications behind items 5 to 54 of Appendix D. The author directed the experiments and the submissions, with the exception stated, and is responsible for every claim.
